\documentclass[a4paper,11pt]{article}

\usepackage{revy}
\usepackage[utf8]{inputenc}
\usepackage[T1]{fontenc}
\usepackage[danish]{babel}


\revyname{MatematikRevyen}
\revyyear{2026}
% HUSK AT OPDATERE VERSIONSNUMMER
\version{0.1}
\eta{$2$ minutter}
\status{Færdig}

\title{Kvantekævle}
\author{Victoria 19', Baldur 20'}

\begin{document}
\maketitle

\begin{roles}
\role{A}[Skuespiller] MFB arrangør + planlægger af kvantekævle
\role{S1}[Skuespiller] Studerende til MFB
\role{S2}[Skuespiller] Studerende til MFB
\role{S3}[Skuespiller] Har en replik
\role{Sn}[Skuespiller] n statist-studerende til MFB
\end{roles}

\begin{props}
\prop{Rekvisit} MFB hehe
\prop{Rekvisit} En tavle/stort stykke papir, hvor man kan signe up til kvantekævle
\prop{Rekvisit} Alting til kævle
\end{props}


\begin{sketch}
\scene{Vi er til MFB og folk står og hygger, snakker og danser. A får en mikrofon stukket i hånden, og stiller sig op på en stol for at give fællesbeskeder}

\says{A} Hej allesammen og velkommen til MFB. I stedet for den sædvanlige beerpong-turnering prøver vi i denne blok noget nyt, nemlig kvantekævle! 

\scene{Folk ser forvirrede ud}

\says{A} Ja, jeg havde heller ikke hørt om det, før jeg crashede fysikkernes Kæmpe Fest. Derfor har jeg allieret mig med S1 og S2, som vil demonstrere.

\scene{S1 og S2 sætter op til kævle på scenen. De resterende studerende stiller sig op omkring og ser på}

\says{A} For at gøre reglerne lidt nemmere at forstå, så kalder vi S1 for Alice og S2 for Bob. 
I kvantekævle kaster vi ikke en sko mod flasken, men en kvanteoperator.

\scene{A tager en gummistøvle frem og giver den til S1}

\says{A} Når Alice så har brugt sin kvanteoperator på flasken 

\scene{S1 kaster gummistøvlen mod flasken og rammer ikke. De studerende som ser på, er skuffede og ærgelige på S1's vegne.}

\says{A}[Mod de studerende] Nej nej, vi ved jo slet ikke om flasken er væltet endnu. Bob skal først \textit{måle} på den, for at finde ud af udkommet.

\scene{S2 løber op og måler på flasken, og konstaterer at den ikke er væltet. S2 tager kvanteoperatoren med tilbage til sin plads.}

\says{A} Arh, bedre held næste gang Alice. Nu går det på tur, ligesom almindelig kævle. 

\scene{S2 gør klar til at kaste, men pludsligt bliver deres øl til en Ice. Alle ser forvirrede ud.}

\says{A} Nårh ja, det glemte jeg at sige. Bobs øl har lavet et bitflip, og er blevet til en Ice. Det kan ske på et vilkårligt tidpsunkt i løbet af spillet. Så, Bob, kast kvanteoperatoren. 

\scene{S2 kaster kvanteoperatoren, og vælter den. S2 begynder at drikke sin Ice}

\says{A}[Til S1] Så, skynd dig nu at måle på flasken Alice! 

\scene{S1 løber ind for at måle på flasken, men S2 har allerede drukket sin Ice. De studerende jubler}

\says{A} Ja, ser det ikke sjovt ud? Og det bedste er, at Alice og Bobs drikkevare er entangled, så når Bob har bundet sin Ice, så har han faktisk også bundet Alices øl! 

\scene{S1 vender sin øl på hovedet, som er tom. S1 og S2 jubler begge som vindere. Måske der bare skal være lys ned her?}

\says{S3}[Til sin sidemakker] Det er simpelthen for mærkligt det her. Jeg går på Caféen?
\scene{Tæppe}

\end{sketch}
\end{document}
